% Deep correctness review of "Algebras Determined by a Braided-Symmetric Product" (v7)
% Build: latexmk -pdf deep-review-v7.tex   (or: pdflatex twice)
\documentclass[11pt,a4paper]{article}

\usepackage[T1]{fontenc}
\usepackage[utf8]{inputenc}
\usepackage{amsmath,amssymb}
\usepackage{booktabs}
\usepackage{array}
\usepackage{url}
\usepackage[margin=2.4cm]{geometry}
\usepackage{enumitem}
\usepackage[hidelinks]{hyperref}
\usepackage[protrusion=true,expansion=false]{microtype}

\newcommand{\yes}{\checkmark}
\newcommand{\no}{$\times$}
\newcommand{\cpt}[1]{\texttt{#1}}          % code / short identifier
\newcommand{\pp}[1]{\path{#1}}             % path: breakable at / - . _
\urlstyle{tt}
\setlength{\emergencystretch}{2.5em}
\newcolumntype{L}[1]{>{\raggedright\arraybackslash}p{#1}}
\setlist[itemize]{leftmargin=1.4em,itemsep=2pt,topsep=3pt}

\title{\bfseries Deep correctness review of\\[2pt]
\emph{Algebras Determined by a Braided-Symmetric Product}, version 7}
\author{Review report}
\date{Repository \cpt{nucleus/work/projects/\ldots/repo}, commit \cpt{977d4ab} (\cpt{braided-sym-zpd v7}).
Compiled \today.}

\begin{document}
\maketitle

\begin{abstract}
\noindent
\textbf{The mathematics holds up with one exception, and the exception is real.} As printed in v7,
Theorem~\cpt{thm:zpd-fxy} is \emph{false} without the unitality hypothesis that its proof uses: a
counterexample is given under D8, and it is the one place in the paper where a stated result fails.
Everywhere else no false statement and no real proof gap was found --- in the criterion, the injective
converse, the codimension-one theorem and the matrix proposition --- and every claim that can be checked
independently was recomputed, including a from-scratch verification of the $\mathbb{F}_3$ computation,
which is the paper's only genuinely \emph{computed} statement. The defects are in the artifacts around the paper: a stale version string
visible in the published PDF, a symbol left dangling by the v7 edit, an abstract still advertising the
v6 formulation, a README that lost all its inline mathematics, a promise of git tags that do not exist,
build instructions for a Lean file that has never elaborated, one wrong citation locator, and one
abstract sentence that is stronger than the theorem behind it. D8 is the one defect that makes a published
statement false, and the working tree adds the missing hypothesis; the rest are packaging, of which
D1--D5 are visible in the published artifacts and should be fixed before the next release. A fix pass
applying D1--D19 is recorded separately (\cpt{results/fixes-applied.md}).
\end{abstract}

\section{Scope}

The review covers the whole repository at commit \cpt{977d4ab}: the LaTeX source
\pp{paper/braided-sym-zpd.tex}, the shipped PDF \pp{pdf/braided-sym-zpd-v7.pdf}, \cpt{README.md}, and
\pp{lean/Research/BraidedSymZpd/BraidedSymZpd.lean}. The paper is short (432 lines of source, 7 pages),
so every proof was read in full rather than sampled.

Three kinds of evidence are distinguished throughout:
(i)~recomputation, where a claim admits a mechanical check;
(ii)~verification against the sources the paper cites;
(iii)~hand checking of the proofs.
Where none of the three was available, this is stated explicitly in Section~\ref{sec:unverifiable} rather
than glossed over.

\section{What was checked mechanically}
\label{sec:mech}

\begin{center}\small
\begin{tabular}{@{}L{0.275\textwidth}L{0.36\textwidth}L{0.29\textwidth}@{}}
\toprule
\textbf{Claim} & \textbf{Method} & \textbf{Result}\\
\midrule
Paper compiles & \cpt{latexmk -pdf} & OK, 7~pp., no undefined references, no overfull boxes\\
\midrule
Shipped PDF is the build of the shipped source & \cpt{pdftotext} diff of shipped vs.\ fresh build
  & identical apart from the \cpt{\textbackslash today} stamp --- in sync\\
\midrule
\cpt{prop:matrix}: exactly 417 zero pairs among $81\cdot81$ & independent Python over $\mathbb{F}_3$
  & \yes\ 417\\
\cpt{prop:matrix}: $\operatorname{rank}\bar\mu_\alpha = 4$, hence $\dim\ker = 12$
  & elimination in Python, plus an analytic argument (the image contains $E_{12}$, $E_{21}$,
    $E_{11}+E_{22}$, $2E_{11}$) & \yes\ 4 and 12\\
\cpt{prop:matrix}: the 417 tensors span a space of dimension 12
  & Python with left-to-right \emph{and} right-to-left pivoting; rank of the transposed matrix; a
    \cpt{sympy} pass over $GF(3)$ & \yes\ 12\\
\cpt{prop:matrix}: conclusion $Z_\mu = \ker\bar\mu$
  & not by dimensions: an explicit 12-vector kernel basis was built and each vector checked to lie in
    $\operatorname{span}(Z)$; with $Z \subseteq \ker$ automatic, this is equality & \yes\ $Z=\ker$\\
\midrule
``$12=12$ for the ordinary product'' & same pipeline & \yes\ 417 / 12 / 12\\
``$13=13$ for the twist by the transpose'' & same pipeline, $\alpha=\mathsf{T}$ as a braiding
  ($xy + y^{\mathsf T}x^{\mathsf T}$) & \yes\ 513 / 13 / 13\\
\midrule
The twisted product ``has no form $f(xy)$'' & exhibit a pair with $xy=0 \ne \mu(x,y)$
  & \yes\ ($x=E_{12}$, $y=E_{11}$)\\
417 is not an artefact of the twisted product & count the pairs with $uv+vu=0$ in $M_2(\mathbb{F}_3)$
  & also 417 \yes\\
\cpt{ex:dual} (dual numbers) & same criterion over $\mathbb{F}_5$: twisted product
  $\operatorname{rank}\bar\mu=1$, $\dim\ker=3$, $\dim\operatorname{span}(Z)=3$ (determined);
  associative product rank~2, $\dim\ker=2$, $\dim\operatorname{span}(Z)=1$ (not zpd)
  & \yes\ both halves hold from the criterion alone\\
\midrule
Lean \cpt{M2F3Check} block (66 lines, core Lean) & \cpt{lean} on the extracted namespace
  & compiles; all three \cpt{\#guard}s pass (417, rank 12, rank $\bar\mu$ 4)\\
Lean file \emph{as a whole} & \cpt{lake build}; then with both missing modules stubbed & \no\ does not elaborate even then: two errors at 124:6 and 233:6 (D5)\\
no \cpt{sorry}, \cpt{admit}, \cpt{axiom}, \cpt{native\_decide} & \cpt{grep} & none found \yes\\
Lean encoding faithful & read \cpt{mul}, \cpt{alpha}, \cpt{mu}
  & \cpt{alpha} is conjugation by $\operatorname{diag}(1,-1)$ \yes; \cpt{mu x y} $=xy+\alpha(yx)$ \yes\\
\bottomrule
\end{tabular}
\end{center}

The Lean \cpt{\#guard}s return exactly the numbers computed here by other means, which is good evidence
that the formal computation is not decorative.

\section{The mathematics, claim by claim}
\label{sec:math}

\begin{itemize}
\item \textbf{Def.~\cpt{def:determined}, Prop.~\cpt{prop:bresar-zpd}, Cor.~\cpt{cor:criterion}.}
The equivalence ``$A$ determined by $\mu$ $\iff$ $Z_\mu=\ker\bar\mu$'' is applied correctly: condition
(iii) of Bre\v{s}ar's Proposition~1.3 specialised to $(A,\mu)$ says exactly $\ker\bar\mu\subseteq Z_\mu$,
and $Z_\mu\subseteq\ker\bar\mu$ is automatic. Neither finite-dimensionality nor surjectivity is needed.
\textbf{Correct.}

\item \textbf{Thm.~\cpt{thm:zpd-fxy}} ($A$ zpd $\Rightarrow$ $A$ determined by $f(xy)$). The one
non-obvious step is $\ker f\subseteq\ker\psi_0$, obtained from $\mu(b,1)=f(b)=0$; then $\psi_0$ factors
through $A/\ker f\cong\operatorname{im} f$ and is extended to $A$. The argument is right and complete,
but it \emph{uses the unit}, which the statement does not mention (D8).

\item \textbf{Prop.~\cpt{prop:injective}.} For injective $f$, $\mu(x,y)=0\iff xy=0$, and then
$\varphi(x,y)=(\psi\circ f)(xy)$. \textbf{Correct.}

\item \textbf{Cor.~\cpt{cor:identity}.} (i) is the theorem with $f=2\operatorname{id}$; (ii) is the
proposition, since $2\operatorname{id}$ is injective iff $\operatorname{char}K\ne2$. The characteristic~2
remark ($\mu\equiv0\Rightarrow$ every admissible $\varphi$ is $0\Rightarrow$ every algebra is determined)
is also correct. \textbf{Correct.}

\item \textbf{Cor.~\cpt{cor:twisted}, Cor.~\cpt{cor:anti}, Cor.~\cpt{cor:triangular}.} In all three the
identity $\alpha(y)\alpha(x)=\alpha(xy)$ is exactly what is needed, and it holds for an endomorphism
(commutative case) and an anti-endomorphism alike; $\alpha$ need not be unital or injective.
\textbf{Correct.}

\item \textbf{Thm.~\cpt{thm:codim}.} The splitting $x=(x-c(x)1)+c(x)1$ gives
$\varphi(x,y)=c(x)c(y)\varphi(1,1)$ and $f(xy)=c(x)c(y)f(1)$; $f(1)\ne0$ supplies a left inverse $g$ of
$a\mapsto a f(1)$. Every hypothesis is used and only assumed ones are used. Consistency check: the
hypotheses force $c(1)=1$, so they are not contradictory. \textbf{Correct.}

\item \textbf{Ex.~\cpt{ex:dual}.} $f=\operatorname{id}+\alpha$ has kernel $K\varepsilon$, a codimension-one
ideal, and $f(1)=2\ne0$; the only idempotents of $K[\varepsilon]/(\varepsilon^2)$ are $0$ and $1$ (from
$a^2=a$, $2ab=b$), so it is not generated by idempotents, and Bre\v{s}ar's Theorem~2.27 gives ``not zpd''.
\textbf{Correct}; the idempotent claim was checked independently, and both halves were also confirmed
numerically over $\mathbb{F}_5$ (Section~\ref{sec:mech}). The ``not zpd'' half moreover admits a direct
proof that avoids the recalled Theorem~2.27 entirely: $\varphi(x,y)=ad$ for $x=a+b\varepsilon$,
$y=c+d\varepsilon$ vanishes on every zero product (if $a=0$ trivially; if $a\ne0$ then $c=0$, and then
$d=0$), while $\varphi(1,\varepsilon)=1\ne0=\varphi(\varepsilon,1)$ although
$1\cdot\varepsilon=\varepsilon\cdot1$, so $\varphi$ does not factor through the product.

\item \textbf{Prop.~\cpt{prop:matrix}.} Verified, see Section~\ref{sec:mech}. The proof as written is a
faithful description of a computation that does what it says.

\item \textbf{Rem.~\cpt{rem:derivations}.} ``$D$ is a braided-symmetric derivation iff $S=D$ on the span
of the products'' follows correctly from the definition of ``determined''.

\item \textbf{Framing claims.} The three symmetrised products are $2xy$, $xy+yx$,
$xy+\alpha(y)\alpha(x)$ --- correct; ``the twisted product describes algebras the zpd theory does not
reach'' --- correct (dual numbers).
\end{itemize}

\noindent
Beyond these hand checks, the following were established exhaustively by independent reviewers and
reproduced here. The criterion $Z_\mu=\ker\bar\mu$ was tested against the \emph{definition} for all
bilinear $\mu$ on $2$-dimensional spaces over $\mathbb{F}_2$ (256) and $\mathbb{F}_3$ (6561) and for 3000
random $\mu$ on $3$-dimensional spaces over both fields, with no mismatch; \cpt{thm:zpd-fxy} and
\cpt{prop:injective} were tested for every unital algebra of dimension $\le3$ over $\mathbb{F}_2$ and
$\le2$ over $\mathbb{F}_3,\mathbb{F}_5$ against every linear $f$ (about 55\,000 pairs, no failure);
\cpt{thm:codim} was tested on 165 (algebra, $f$) cases satisfying its hypotheses, including a
noncommutative instance in $\mathrm{Tri}(M_2(\mathbb{F}_p))$ where
$Z_\mu=N\otimes A+A\otimes N=\ker\bar\mu$ holds with dimension $8=8$; the recalled Theorems~2.27 and 7.2
were checked on the small unital algebras; and the content of Corollary~2.17(b) can be proved outright:
$\left[\begin{smallmatrix}1&m\\0&0\end{smallmatrix}\right]$ is idempotent because $m\cdot1=m$, and
$\left[\begin{smallmatrix}0&m\\0&0\end{smallmatrix}\right]$ is its difference with
$\left[\begin{smallmatrix}1&0\\0&0\end{smallmatrix}\right]$, so $M$ lies in the span of the idempotents
and $\mathrm{Tri}(A_1,M,A_2)$ is generated by them.

\section{Defects}
\label{sec:defects}

Severity: D8 is a false statement in v7 as printed and comes first in importance; D1--D5 are visible in
the published artifacts and should be fixed before the next release; D6, D7 and D18 are referee-visible;
D9--D11 are minor; D12--D17 and D19 are cosmetic, provenance, or a missing one-line justification.

\subsection*{D1. The published PDF says ``Preprint version 2'', everything else says v7}
\cpt{paper/braided-sym-zpd.tex:53} has \cpt{\textbackslash newcommand\{\textbackslash paperversion\}\{2\}},
which prints ``Preprint version 2, compiled 2026/10/17'' on the title page of
\cpt{pdf/braided-sym-zpd-v7.pdf}, whose file name, the README and the commit all say v7. The counter was
never updated after v2. \emph{Fix:} \cpt{\textbackslash paperversion\{7\}}, or drop the macro.

\subsection*{D2. \cpt{q\_\textbackslash circ} is undefined --- a regression introduced by the v7 edit}
Line 290: ``When $\operatorname{char}K\ne2$ the map $q_\circ$ is onto, because
$x\circ\tfrac12=x$.'' In v6 the companion $q_\mu\colon (A\otimes A)/Z_\mu\to A$ was defined (line 157)
and used throughout the proof of the criterion. v7 replaced that proof by the citation of
Bre\v{s}ar's Proposition~1.3 and deleted the defining sentence, but left this one use. With
Proposition~1.3 doing the work, the sentence is also redundant: the criterion needs no surjectivity.
\emph{Fix:} delete the sentence, or define $q_\circ\colon A\otimes A\to A$, $x\otimes y\mapsto x\circ y$.

\subsection*{D3. \cpt{README.md} has lost all inline mathematics}
The file contains zero \cpt{\$} characters. Three visible damages in the summary paragraph (line 8):
``\ldots on a commutative algebra, \textbf{has the form} $f(xy)$'' (the subject $\mu$ is gone),
``determined by every such \textbf{.}'' ($\mu$ gone), ``the algebra of \textbf{2 2} matrices'' (the
$\times$ is gone). \emph{Fix:} restore the three expressions.

\subsection*{D4. README promises git tags that do not exist}
``Each version is also a tag of this repository'' --- \cpt{git tag} is empty, and the history begins at
v3 (seven commits; v1 and v2 are absent). Either create the tags or reword.

\subsection*{D5. \cpt{cd lean \&\& lake build} does not work, and cannot be made reproducible as pinned}
Four independent reasons:
\begin{enumerate}[label=(\arabic*),leftmargin=1.8em,itemsep=2pt]
\item the file imports \pp{Research.Zpd.Basic} and \pp{Research.Zpd.DualNumbers}, which are \emph{not
in the repository}: 479 of the 545 lines depend on them, including
\pp{Research.Zpd.ZeroProductDetermined} and \pp{not_zeroProductDetermined_dualNumbers}. The README's
description ``Lean 4 files that check computations and lemmas of the paper (mathlib only)'' (line 24) is
therefore false as well;
\item \textbf{even with both modules stubbed, the file does not elaborate.} It fails at 124:6 and 233:6
with \emph{Alternative \cpt{zero} has not been provided}, because the two \cpt{induction} blocks over
$B\otimes_k B$ in \cpt{mapDetermined\_iff\_injective} and \cpt{mapDetermined\_iff\_injective$'$} omit the
\cpt{zero} branch. This was reproduced against the installed Mathlib (Lean 4.33.1): a minimal
\cpt{induction s with | tmul \ldots\ | add \ldots} on a tensor fails with exactly that message, and the
same proof compiles once a \cpt{| zero =>} branch is added. So the artifact has never been elaborated,
and its failure is not only the missing dependency;
\item \pp{lakefile.toml} pins mathlib to \cpt{rev = "master"} --- not a reproducible pin --- and
current master requires \cpt{lean4:v4.36.0-rc2}, while \pp{lean-toolchain} pins \cpt{v4.35.0}. A build
therefore \emph{rewrites the tracked} \pp{lean-toolchain} (it did, on the review run; it was restored)
and pulls 625\,MB;
\item no \pp{lake-manifest.json} is committed.
\end{enumerate}
\emph{Fix:} vendor the two \pp{Research/Zpd} modules, add the \cpt{zero} branches, pin mathlib to a tag
or commit, commit the manifest, align \pp{lean-toolchain}. Until then the README should not claim the
file builds.

\subsection*{D6. Citation locator: \cpt{Tri}$(A_1,M,A_2)$ is cited to \S2.1, which is ``zpd Rings''}
Line 321 cites \cpt{[\S2.1]} \emph{twice in a row} for the notation $\mathrm{Tri}(A_1,M,A_2)$. The
book's \S2.1 is ``zpd Rings''; the construction is an example of a zpd algebra, i.e.\ \S2.2 (``Examples
and Non-examples of zpd Algebras''), which Corollary~2.17(b) then uses. \emph{Fix:} \cpt{[\S2.2]}, once.

\subsection*{D7. The abstract (and README) is stronger than Cor.~\cpt{cor:twisted}}
``For the symmetrised products of three braidings on a commutative algebra, $\mu$ has the form $f(xy)$''
is \emph{false for a merely linear} $\alpha$, which is what the Introduction allows when it defines the
twisted swap. Take $A=K[\varepsilon]/(\varepsilon^2)$ (commutative) and the linear involution
$\alpha(a+b\varepsilon)=b+a\varepsilon$; then $\alpha(\varepsilon)=1$ and
\[
\mu(\varepsilon,\varepsilon)=\varepsilon^2+\alpha(\varepsilon)\alpha(\varepsilon)=1,
\qquad \varepsilon\cdot\varepsilon=0 .
\]
If $\mu(x,y)=f(xy)$ with $f$ linear, then $f(0)=\mu(\varepsilon,\varepsilon)=1$, contradicting
$f(0)=0$. \S4 states the result correctly, with ``$\alpha$ an algebra endomorphism''.
The repair is exactly sharp: if moreover $\alpha(1)=1$ then $\mu_\alpha$ is of the form $f(xy)$
\emph{if and only if} $\alpha$ is multiplicative (put $x=1$ to get $f(y)=y+\alpha(y)$ and compare with
$\mu_\alpha(x,y)$), so ``endomorphism'' is not merely sufficient but the exact missing hypothesis.
\emph{Fix:} add ``for $\alpha$ an endomorphism'' in the abstract and the README.

\subsection*{D8. \cpt{thm:zpd-fxy} is false as stated: the unit is needed}
The theorem says ``Let $A$ be zpd, let $f\colon A\to A$ be linear'' and its proof uses $\mu(b,1)=f(b)$
for $b\in\ker f$. Because \S3's preamble assumes $A$ unital, and the abstract and the introduction both
say ``unital'', this was first recorded here as statement hygiene only, with the remark that the
hypothesis could be dropped rather than added: $\psi_0$ vanishes on $A^2\cap\ker f$, so extend
$\psi_0|_{A^2}$ to $A$ killing $\ker f$. \textbf{That reasoning is wrong, and the theorem is false
without unitality.} From $b=\sum_i x_i y_i$ and $f(b)=\sum_i\mu(x_i,y_i)=0$ one cannot conclude
$\sum_i\varphi(x_i,y_i)=0$, and applying Bre\v{s}ar's condition (ii) to $\mu$ would assume what is being
proved.

\emph{Counterexample} (found by a reviewer of this report; recomputed here over $\mathbb{F}_2$,
$\mathbb{F}_3$, $\mathbb{F}_5$, $\mathbb{F}_7$). Let $A=\langle e_0,e_1,e_2\rangle$ with
$e_0^2=e_0$, $e_0e_1=e_0e_2=e_2$, $e_1e_0=e_1-e_2$ and all other products zero. Then $A$ is
associative, $A^2=A$, $A$ has no unit, and $A$ \emph{is} zpd: $\dim Z_m=\dim\ker\bar m=6$. Take
$f(a,b,c)=(a,0,c)$, so that $\ker f=\langle e_1\rangle$, and let $\mu(x,y)=f(xy)$. Then
$\dim Z_\mu=6$ but $\dim\ker\bar\mu=7$, so by Corollary~\cpt{cor:criterion} $A$ is \emph{not} determined
by $\mu$. Explicitly, $\varphi(x,y)=x_0(y_1+y_2)$ vanishes on every $\mu$-zero pair, while
$\varphi(e_0,e_1)=1$ and $\varphi(e_1,-e_0+e_2)=0$ although $\mu(e_0,e_1)=\mu(e_1,-e_0+e_2)=e_2$.

\emph{Fix (applied):} the statement now reads ``Let $A$ be a unital zpd algebra'', which is what the
abstract and the introduction already said. The hypothesis must be \emph{added}, not dropped.

\subsection*{D9. Two change logs for v7 that describe different changes}
The source header says v7 is the reworking of the introduction around Bre\v{s}ar's Theorem~7.2, plus
\cpt{rem:preserving}, plus ``claims added: \cpt{cor:triangular}, \cpt{prop:matrix}''. The README's v7
entry says v7 is the retitling towards \cpt{braided-sym-deriv}, the criterion now cited from
Proposition~1.3, and \cpt{rem:derivations} --- and does not mention the two new claims. Both edits are in
the same commit; each log describes half of it.

\subsection*{D10. Stale or incorrect statements in the Lean file}
\begin{itemize}[leftmargin=1.4em,itemsep=1pt]
\item Line 45 says the noncommutative case ``(for example, $M_2(\mathbb{F}_3)$ with the transpose) is
open'', but the same file then checks $M_2(\mathbb{F}_3)$, and the \emph{transpose} case is settled
affirmatively by \cpt{cor:anti}, since the transpose reverses products. Wrong twice.
\item The module docstring advertises \cpt{symProdLinear} and
\cpt{mapDetermined\_symProd\_std\_iff\_zJpd} as the file's main definitions and results; neither is
defined anywhere in it (both occur only in the docstring, lines 26, 35 and 102).
\item \cpt{cor:twisted} is formalised as \cpt{mapDetermined\_twisted\_of\_comm\_zpd} for
$\alpha\colon B\to_a[k]B$, i.e.\ a \emph{unital} endomorphism, whereas the paper's corollary only needs
$\alpha$ multiplicative (no theorem of the paper assumes $\alpha(1)=1$). The formal statement is
therefore strictly weaker than the published one.
\item Lines 153 and 172 refer to \cpt{BraidedSymmetric.lean}, and line 45 to
\cpt{formalization-debt.md}; neither file is in this repository.
\end{itemize}

\subsection*{D11. Ambiguity in ``the twist by the transpose'' (line 413)}
The paper writes the twisted product both as $xy+\alpha(y)\alpha(x)$ (the braiding) and as
$xy+\alpha(yx)$ (\cpt{prop:matrix}); these agree for automorphisms but not for anti-automorphisms. For
$\alpha={}^{\mathsf T}$ the two readings give $13=13$ (the braiding, as printed \yes) and $12=12$
($xy+\alpha(yx)$, also determined, computed here). The printed number is right for the intended reading;
the notation should say which product is meant.

\subsection*{D12. The abstract advertises a ``quotient criterion'' the v7 body no longer states}
Line 72: ``The answer is a quotient criterion that needs no surjectivity hypothesis.'' In v7 the
criterion is the kernel equality $Z_\mu=\ker\bar\mu$; the words \emph{quotient} and \emph{surjective}
occur nowhere else in the body except inside the recalled Theorem~7.2. Same root cause as D2: v6 stated
the criterion as ``$q_\mu$ is injective'', with a left inverse replacing an inverse and no surjectivity
assumed, and v7 deleted that formulation without touching the abstract. Low severity --- the kernel
equality is equivalent to injectivity on $(A\otimes A)/Z_\mu$ --- but the abstract describes a
presentation the paper does not give.

\subsection*{D13. The abstract names an object that does not exist}
Line 71: ``For the symmetrised products of braided-symmetric derivations, this asks when the algebra with
that product is zero product determined.'' A symmetrised product is attached to a \emph{braiding}; a
braided-symmetric derivation is a \emph{derivation of} $(A,\cdot_\sigma)$. There is no ``symmetrised
product of a derivation'' anywhere in the paper. Relatedly, the title and keywords say
``braided-symmetric product'', a term that occurs in no line of the body (which says ``symmetrised
product'' and ``braided-symmetric derivation''). Added in v7.

\subsection*{D14. \S2 restates the symmetrised product with $\mu$ for the algebra product}
Line 173: ``the symmetrised product $x\cdot_\sigma y=\mu(x\otimes y)+\mu(\sigma(x\otimes y))$, where
$\mu$ is the product of $A$''. But $\mu$ is the arbitrary bilinear map of \S\S3--5 (line 199), the product
is $m$ (line 108), and $\mu$ is defined on $A\times A$, not on $A\otimes A$. The formula also disagrees
\emph{in notation} with the Introduction's $x\cdot_\sigma y=m(x\otimes y+\sigma(x\otimes y))$ (line 117);
the two agree mathematically, so this is a notational clash rather than an error, but it sits
immediately before \cpt{def:adrabi-bsder}, which therefore has two incompatible definitions of
$\cdot_\sigma$ around it.

\subsection*{D15. The v7 header comment misattributes one change and omits others}
Line 1 lists ``Acknowledgments removed'' as a v7 change, but the section was deleted one commit earlier,
in v6 (\cpt{git diff 81395eb e56e0a7}). The header also omits two real v7 edits: the abstract clause ``as
the transpose does on matrices'' was deleted, and ``$M_n(K)$ is determined by $xy+y^{\mathsf T}x^{\mathsf
T}$'' was replaced by the more general statement that \cpt{cor:anti} applies ``to $M_n(K)$ and to the
transpose, and to every linear map $\alpha$ of $M_n(K)$ that reverses products''.

\subsection*{D16. README changelog order}
It announces ``newest first'' but lists v1 (2026-09-25) before v2 (2026-09-27); and neither v1 nor v2
exists anywhere in the repository history (see D4).

\subsection*{D17. Cosmetic}
Duplicate citation of one locator (\cpt{[\S2.1] \ldots [Section 2.1]}, line 321); three labels are never
\cpt{\textbackslash ref}'d (\cpt{cor:identity}, \cpt{def:bresar-zpd}, \cpt{def:bresar-zpp}) --- harmless.

\subsection*{D18. The Introduction credits matrix algebras to Cor.~\cpt{cor:triangular}}
Line 129: ``So the zpd algebras of Bre\v{s}ar, such as matrix algebras and triangular algebras
[Corollary~2.17], are determined by all these products (Corollary~\cpt{cor:triangular}).'' But
\cpt{cor:triangular} covers only $\mathrm{Tri}(A_1,M,A_2)$; $M_n(B)$ is not of that form, and is reached
by \cpt{thm:zpd-fxy} together with Corollary~2.17(a). As written the sentence reads as a claim about both
families. \emph{Fix:} cite \cpt{thm:zpd-fxy} (and Cor.~2.17) for matrix algebras and \cpt{cor:triangular}
for the triangular ones.

\subsection*{D19. Two unproved assertions (both true, one line each)}
Line 395: ``When $\alpha$ is an automorphism of a noncommutative algebra, the twisted product
$xy+\alpha(yx)$ is not of the form $f(xy)$, and Theorems~\cpt{thm:zpd-fxy} and \cpt{thm:codim} do not
apply''; and line 292: the Jordan product ``is not of the form $f(xy)$ on a noncommutative algebra''.
Both are true --- $y=1$ forces $f=\operatorname{id}+\alpha$ respectively $f=2\operatorname{id}$, and then
$\alpha(yx)=\alpha(xy)$, respectively $xy+yx=2xy$, gives $yx=xy$, contradicting noncommutativity in
characteristic $\ne2$ --- but the text offers no argument for either, and neither sentence carries a
\cpt{claim-ok} marker. \emph{Fix:} one line of proof, or a parenthetical.

\section{Verification against the primary source}
\label{sec:primary}

\begin{center}\small
\begin{tabular}{@{}L{0.40\textwidth}L{0.51\textwidth}@{}}
\toprule
\textbf{Locator in the paper} & \textbf{Status}\\
\midrule
Definition 1.1 (zpd, nonassociative) & \textbf{verbatim match} \yes\\
Proposition 1.3 (i)--(iv) & \textbf{verbatim match} \yes\ (other authors cite it as \cpt{[3, Proposition 1.3]}
  for the tensor form)\\
Corollary 2.17(a) ($M_n(B)$ is zpd) & \yes\ (cited elsewhere as \cpt{[3, Corollary 2.17 (a)]})\\
Theorem 2.27 (finite-dimensional unital zpd $\iff$ generated by idempotents) & \yes\ --- and not a slip
  for Theorem 3.7: 3.7 is the numbering in Bre\v{s}ar's \emph{Expo.\ Math.}\ 34 (2016) paper, of which
  the book's Theorem 2.27 is the book version (one paper cites the book as \cpt{[5, Theorem 2.27]} and
  the article as \cpt{[6, Theorem 3.7]})\\
\S2.1 for $\mathrm{Tri}(A_1,M,A_2)$ & \no\ \S2.1 is ``zpd Rings'' --- should be \S2.2 (D6)\\
Example 2.19 ($T_n(F)$), Example 2.23 (adjoining a unity) & \textbf{content confirmed} against the
  published \emph{Expo.\ Math.}\ 34 (2016) paper, where the same two items appear as Example~2.19
  ($T_n(F)$ is zpd) and as the propositions stating that adjoining a unity to an algebra with trivial
  multiplication is not zpd; the book's own locator numbers are \textbf{not confirmed}\\
Def.~3.14 (zJpd), Def.~7.1/Thm.~7.2, Ex.~1.10 & plausible, consistent with the book's section layout and
  with the parallel numbering of the \emph{Expo.\ Math.}\ paper, but \textbf{not confirmed} against the
  book itself\\
Horn--Johnson 0.2.5 (``the transpose reverses products'') & the statement is true; the \textbf{locator was
  not confirmed}\\
\bottomrule
\end{tabular}
\end{center}

\section{Claims not verifiable here}
\label{sec:unverifiable}

\noindent\textbf{Access, reference by reference.} Bre\v{s}ar's book (Birkh\"auser 2021) was \emph{not}
directly accessible: the publisher's chapter pages redirect to an institutional login, which was not
worked around. What was used instead is (a) the text of Chapter~1 --- Definition~1.1 and
Proposition~1.3 (i)--(iv) --- as reproduced verbatim by a third-party mirror surfaced by search,
(b) the book's table of contents for the section structure (which is how the \S2.1 error of D6 was
found), and (c) for the remaining locators, only the way other papers cite the book. Bre\v{s}ar's
\emph{Expo.\ Math.}\ 34 (2016) paper, which contains the article versions of Theorem~2.27 and
Corollary~2.17, was available in full from the author's page
(\url{https://users.fmf.uni-lj.si/bresar/documents/zpdidemp.pdf}) and was used to confirm the
\emph{content} of Examples~2.19 and 2.23 and of Theorem~3.7/2.27. Horn--Johnson was available only as
the publisher's free front matter, whose contents list sections but not subsections, so the locator
\S0.2.5 stays unconfirmed (the statement it supports, that the transpose reverses products, is
standard). Adrabi's 2026 preprint is available nowhere: searches of arXiv, the author's public
profiles, and his 2024 doctoral thesis found no trace of the title, so the paper's framing cannot be
checked against it at all.

\begin{itemize}
\item \textbf{Adrabi 2026} (the author's own preprint) --- Definition~3.1 (braided-symmetric derivations),
Remark~3.4, Section~1 and Section~4 are cited for the paper's entire framing, but the preprint is neither
in this repository nor publicly reachable. A referee cannot check them. Consider making it available, or
restating the two facts: they are one-line computations ($D(x\cdot_\sigma y)=D(x)\cdot_\sigma y+x\cdot_\sigma
D(y)$ is literally the definition of a derivation of $(A,\cdot_\sigma)$).
\item The Lean file's 479 non-computational lines: unverifiable until D5 is fixed. What can be said is
that the statements they claim to prove are the paper's theorems, and that those were verified by hand.
\end{itemize}

\section{Substantive remarks (not defects)}
\label{sec:remarks}

\begin{enumerate}[label=\arabic*.,leftmargin=1.8em,itemsep=3pt]
\item \textbf{Proposition~\cpt{prop:matrix} may be a shadow of a general theorem.} With
$\alpha$ the conjugation by an involutive $D$, the algebra splits by $D$-parity and on homogeneous parts
$\mu(x,y)=xy+(-1)^{|x|+|y|}yx$. That suggests the $\mathbb{F}_3$, $n=2$ computation could become a
statement for every $n$ and every field of characteristic $\ne2$, which would also sharpen
Question~\cpt{q:noncomm}. As it stands the proposition is a single data point over the smallest field.

\item \textbf{\cpt{thm:codim} and \cpt{cor:twisted} overlap.} On a commutative algebra
$\mu_\alpha=f(xy)$ with $f=\operatorname{id}+\alpha$ and $\ker f=\ker c$ for the augmentation $c$, so the
dual-numbers instance is \cpt{thm:codim} applied to the twisted product. Fine, but worth a linking
sentence.

\item The paper is candid that the criterion is Bre\v{s}ar's Proposition~1.3 re-read; the genuinely new
content is the $f(xy)$ class, the two ``beyond zpd'' classes, and the computation. That is a coherent but
modest contribution, appropriate for a short note --- and the abstract should not oversell it (D7).
\end{enumerate}

\section{Reproduction}

\begin{verbatim}
cd paper && latexmk -pdf braided-sym-zpd.tex          # 7 pp.

sed -n '/^namespace M2F3Check/,/^end M2F3Check/p' \
  ../lean/Research/BraidedSymZpd/BraidedSymZpd.lean > M2F3Check.lean
lean M2F3Check.lean                                   # 3 #guards pass

cd ../lean && lake build      # FAILS: Research.Zpd.Basic, .DualNumbers
git tag                       # empty (README claims tags)
grep -n paperversion ../paper/braided-sym-zpd.tex   # 2, not 7

python3 ../results/m2-verify.py   # all F_3 numbers, with assertions
\end{verbatim}

\noindent
The independent $\mathbb{F}_3$ verification (zero-pair counts, ranks, and the kernel-basis-in-span
argument) is the script \cpt{results/m2-verify.py}; it asserts $417/12/12$, $417/12/12$ and $513/13/13$
for the twisted, ordinary and transpose cases respectively.

\medskip
\noindent\emph{Note on method.} The proofs and the computational claims were checked directly. Seven
directions were carried out by independent reviewers --- version history, the Lean artifact, the
document-level audit, Section 2, Sections 3 and 5, the primary-source audit, and an exhaustive
test-bench over small algebras --- and every finding included here was re-checked against the
repository or recomputed. Four reviewer claims were not adopted as stated: a characterisation of the v7 edit as
weakening a claim about $M_n(K)$ (it generalises it), a non-unital counterexample to line 292
(the Jordan product on strictly upper triangular $3\times3$ matrices is \emph{not} a function
of $xy$, so the example refutes nothing), a claim that the unit of \cpt{thm:zpd-fxy} is
dispensable (false, and this review had itself endorsed it in an earlier draft --- see D8),
and an understated reading of D8 as statement hygiene, which a reviewer's counterexample
corrected.

\end{document}
